% Separate post-CP24 proof. Complete author, root, C and D audits PASS; actual arrows between already realized families.
\subsection{Intrinsic finite-cycle lifting of canonical Hom maps}

\begin{proposition}[Proved: fixed-loss lifting of intrinsic Hom jets]
\label{D:intrinsic-Hom-fixed-loss-jets}
Let $\pi:Y\to\operatorname{Spec}A$ be a projective regular
resolution of a complete normal local surface over an
algebraically closed coefficient field $k$. Let $C$ be a
nonzero effective exceptional Cartier divisor whose support is
the entire closed fibre, and suppose
two elements $x,y$ generate its invertible ideal:
\[
                 (x,y)\mathcal O_Y=\mathcal O_Y(-C).
                                                        \tag{D.HJ.1}
\]
Let $J$ be an ideal of the fixed base ring acting on $A$, with
$x,y\in J$ and $JH^1(Y,\mathcal O_Y)=0$.
Let $T$ be a finite-dimensional local Artinian $k$-algebra with
residue field $k$, and let $F_T,G_T$ be vector bundles on $Y_T$.
Assume that, for fixed positive integers $a,b$, each of
$x^a\operatorname{Id}_{F_T}$ and $y^a\operatorname{Id}_{F_T}$
factors through a finite free $\mathcal O_{Y_T}$-module,
and each of $x^b\operatorname{Id}_{G_T}$ and
$y^b\operatorname{Id}_{G_T}$ does also.

Put $\ell=a+b+1$ and
$H_T=\mathcal Hom(F_T,G_T)$. For every integer $d$ and every
$k\ge\max\{0,d\}$, the inclusion induces the zero map
\[
 H^1\bigl(Y_T,H_T((d-k-\ell)C)\bigr)
       \longrightarrow H^1\bigl(Y_T,H_T((d-k)C)\bigr).
                                                        \tag{D.HJ.2}
\]
Consequently, for $k\ge\max\{1,d\}$,
\[
 \begin{split}
 &\operatorname{im}\bigl[
 H^0((k+\ell)C_T,H_T(dC)|_{(k+\ell)C_T})
       \longrightarrow H^0(kC_T,H_T(dC)|_{kC_T})\bigr]\\
 &\hspace{12mm}=
 \operatorname{im}\bigl[
 H^0(Y_T,H_T(dC))
       \longrightarrow H^0(kC_T,H_T(dC)|_{kC_T})\bigr].
 \end{split}                                           \tag{D.HJ.3}
\]
Thus a map on the higher cycle has a global map with precisely
its prescribed lower-cycle restriction. The assertion holds
after every Artinian coefficient change, using the pulled-back
factorizations. It does not assert that either raw image
commutes with coefficient change.

When $a=b=1$ and $d=0$, the loss is three. If a given higher-cycle
map is an isomorphism, its global lift matching the lower cycle
is an isomorphism on $Y_T$.
\end{proposition}

\begin{proof}
At a local trivialization of $\mathcal O_Y(-C)$, write
$x=su$, $y=sv$ with $(u,v)=1$. Then $(u^j,v^j)=1$ for every
$j\ge1$: otherwise a prime would contain both $u$ and $v$.
Thus $x^j,y^j$ generate $\mathcal O_Y(-jC)$ and give the
locally split pencil sequence
\[
 0\longrightarrow\mathcal O_Y(jC)
 \longrightarrow\mathcal O_Y^{\,2}
 \xrightarrow{(x^j,y^j)}\mathcal O_Y(-jC)
 \longrightarrow0.                                    \tag{D.HJ.4}
\]
All coherent cohomology on $Y_T$ in degrees at least two
vanishes. Indeed properness gives finite cohomology modules
over the complete local base $A\otimes_kT$; formal functions
compute their completions on thickenings of the one-dimensional
closed fibre, where those degrees vanish. This is the same
exact fibre-dimension argument used in the relative realization
and Cartier-jet theorems.

The scalar base-change equality over the field $k$ shows
$JH^1(Y_T,\mathcal O_{Y_T})=0$. Apply (D.HJ.4) after base
change and use $H^2(\mathcal O_{Y_T}(jC))=0$. It gives a
surjection from two copies of $H^1(\mathcal O_{Y_T})$ onto
$H^1(\mathcal O_{Y_T}(-jC))$, so
\[
       JH^1(Y_T,\mathcal O_{Y_T}(-jC))=0
                              \quad(j\ge0).             \tag{D.HJ.5}
\]

Dualizing the assumed factorization on $F_T$ gives a
factorization of $x^a\operatorname{Id}_{F_T^\vee}$ through
a free bundle. Tensoring it with the factorization on $G_T$
therefore factors $x^{a+b}\operatorname{Id}_{H_T}$ through
a free bundle. The same holds for $y^{a+b}$. After tensoring
by $\mathcal O_{Y_T}((d-k)C)$, the induced cohomology maps
factor through sums of
$H^1(\mathcal O_{Y_T}((d-k)C))$. Since $k\ge d$, (D.HJ.5)
annihilates these groups by both $x$ and $y$. Hence
\[
 x^\ell H^1(Y_T,H_T((d-k)C))
 =y^\ell H^1(Y_T,H_T((d-k)C))=0.                         \tag{D.HJ.6}
\]

Tensor the $j=\ell$ sequence (D.HJ.4) by
$H_T((d-k)C)$. Its degree-one cohomology gives a surjection
\[
 H^1(Y_T,H_T((d-k)C))^{\,2}
       \longrightarrow H^1(Y_T,H_T((d-k-\ell)C)).
                                                        \tag{D.HJ.7}
\]
Composing (D.HJ.7) with the inclusion to the shallower
cohomology is multiplication by $x^\ell,y^\ell$ on its two
summands. Both vanish by (D.HJ.6), proving (D.HJ.2).

Now compare the restriction exact sequences for $kC_T$ and
$(k+\ell)C_T$ applied to the vector bundle $H_T(dC)$.
The connecting homomorphism of a section restricted from
the larger cycle is the image of its larger-cycle connecting
class under (D.HJ.2), hence is zero. The smaller-cycle exact
sequence then gives a global lift with that precise restriction.
Conversely every global section restricts through the larger
cycle. This proves (D.HJ.3).

All arguments use actual vector bundles, actual factorizations,
and exact sequences over $T$. Pulling them back to another
allowed Artinian algebra reproduces the proof; no cohomology
base-change inference about the displayed images is made.
Finally, if $d=0$ and the lower-cycle map is an isomorphism,
the determinant of its global lift is invertible at every
point of the closed fibre, since $C$ has full exceptional
support. Its noninvertible locus is closed and proper over
the local base and misses that fibre, so is empty. The
global map is therefore an isomorphism.
\end{proof}

\begin{corollary}[Proved: canonical and Frobenius arrow windows on arbitrary covers]
\label{D:hereditary-canonical-Hom-arrow-window}
Use the fixed Noether-different and Cartier data for the
finite punctured covers $A/R$ of the preceding hereditary
budget theorems. Let $F_T,G_T$ be the canonical full pullbacks
of already realized relative maximal Cohen--Macaulay modules,
for example the families satisfying the retained hypotheses
of Proposition~\ref{D:hereditary-weak-cycle-relative-realization}.
Then (D.HJ.3) holds with loss three for their ordinary Hom
maps. If their closed fibres are the ordinary-unit bundles
in the preceding determinant theorem and the families carry
its fixed-determinant quotient presentations at the retained
levels, then at each fixed pair $k,k+3$ the actual finite-cycle
Hom data have the $O_k(tr_Fr_G)$ size of the quotient-Hom
theorem. No such size assertion is made for arbitrary
maximal Cohen--Macaulay closed fibres.

For the one-step Frobenius Hom
$\mathcal Hom(F_{Y_T}^*F_T,G_T)$, the loss is $p+2$.
For the opposite bounded-pole Hom
$\mathcal Hom(G_T,F_{Y_T}^*F_T)(pC)$, the same loss applies
for $k\ge p$. Here $F_{Y_T}$ is the absolute $p$-Frobenius,
including its action on the Artinian coefficient ring.
These are actual prescribed lower-jet lifts of the indicated
arrows. No unit-compatibility equation or arbitrary object
realization is asserted by this corollary.
\end{corollary}

\begin{proof}
Proposition~\ref{D:cover-uniform-Noether-different} and its
relative source factorization give factorizations of
$x\operatorname{Id}$ and $y\operatorname{Id}$ on both
realized modules through finite free $A_T$-modules. Pull
back and apply double dual. Functoriality gives the same
factorizations on their canonical full bundles through
finite free $\mathcal O_{Y_T}$-modules. The fixed ideal
annihilates scalar $H^1(\mathcal O_Y)$ by
Proposition~\ref{D:cover-uniform-canonical-defect}.
The proposition therefore applies with $a=b=1$.
Its finite-cycle data are the actual Hom spaces computed
by Proposition~\ref{D:hereditary-ordinary-unit-finite-cycle-Hom};
the source modules of the factorization are not extra
matrix-valued coordinates in that presentation.

Absolute Frobenius pulls the first factorization back to
one for $x^p\operatorname{Id}_{F_{Y_T}^*F_T}$, and similarly
for $y^p$. Pullback preserves the free middle term and
the equality of the composite, regardless of flatness of
Frobenius on $T$. Thus the first Frobenius Hom has exponents
$a=p,b=1$, and the opposite one has $a=1,b=p$.
The twist $pC$ in the latter is exactly $d=p$ in
(D.HJ.2)--(D.HJ.3), explaining the condition $k\ge p$.
This proves the claimed uniform arrow windows without
assuming that their Frobenius or composition relations
lift automatically.
\end{proof}

\begin{remark}
The factorizations are available for already realized
canonical families. The proposition does not impose them
on arbitrary finite-cycle bundles, and does not remove
the coflat or split-boundary realization conditions.
Likewise a lift of an individual arrow is not asserted
to preserve independently lifted relations. These
distinctions are necessary in any later unit-cycle
representation count.
\end{remark}
